% TOP_PROBLEM: {"id": 2305052, "title": "Research Problems in Function Theory — Problem 5.52", "category_id": 8, "category": {"id": 8, "name": "analysis", "display_name": "Analysis", "description": "Limits, continuity, calculus, and function theory.", "slug": "analysis", "order_index": 8, "created_at": "2026-07-31T15:26:25.671Z"}, "status": "open", "problem_number": "AMR-022-5052", "set_id": 13}
% FIRST_POSTED: 2026-10-02
% ENTRY_KIND: statement_only
% UnsolvedMath ID 2305052; source catalogue code AMR-022-5052
% !TeX program = lualatex
% Definition and sources only; this file is not an LLM solution attempt.
\documentclass[11pt,a4paper]{article}
\usepackage[margin=25mm]{geometry}
\usepackage{fontspec}
\setmainfont{lmroman10-regular.otf}[BoldFont=lmroman10-bold.otf,
 ItalicFont=lmroman10-italic.otf,BoldItalicFont=lmroman10-bolditalic.otf]
\usepackage{amsmath,amssymb,mathtools}
\usepackage{unicode-math}
\setmathfont{latinmodern-math.otf}
\AtBeginDocument{\let\setminus\smallsetminus}
\usepackage{xurl,hyperref}
\hypersetup{colorlinks=true,urlcolor=blue}
\setcounter{secnumdepth}{0}
\setlength{\parindent}{0pt}
\setlength{\parskip}{5pt}
\setlength{\emergencystretch}{3em}
\begin{document}
\section*{Research Problems in Function Theory — Problem 5.52}\label{2305052}
{\small UnsolvedMath ID 2305052 · AMR-022-5052 · Analysis · AMR Open Problem Lists}
\subsection{Definitions and mathematical statement}
Let $\Gamma$ be a Fuchsian group, meaning a discrete subgroup of the holomorphic automorphism group of $\mathbb D$. It acts also on $\mathbb T=\partial\mathbb D$. Suppose a Borel set $B\subset\mathbb T$ has pairwise disjoint translates in the sense
\[
B\cap\gamma(B)=\varnothing\qquad
(\gamma\in\Gamma\setminus\{\mathrm{id}\}).
\]
For compact $K\subset\mathbb C$, its logarithmic capacity is
$\operatorname{cap}K=\exp[-\inf_\mu\iint\log(1/|z-w|)\,d\mu(z)d\mu(w)]$, where $\mu$ ranges over probability measures supported on $K$. Positive capacity of $B$ means some compact subset has positive capacity.
Is $\operatorname{cap}B>0$ sufficient to force $\Gamma$ to be of convergence type, namely
\[
\sum_{\gamma\in\Gamma}\bigl(1-|\gamma(0)|\bigr)<\infty?
\]
Equivalently one may use $\sum_\gamma e^{-d_{\mathbb D}(0,\gamma(0))}$ with the usual curvature-minus-one hyperbolic distance. Finite stabilizers do not affect convergence. The source gives the converse: convergence-type groups admit such a positive-capacity set. The requested implication therefore asks whether the existence of a sufficiently large disjoint-translates boundary set characterizes convergence type.
\subsection{Short English statement}
Does a boundary set of positive logarithmic capacity with disjoint Fuchsian-group translates force convergence of the group's Poincaré series?
\subsection{Sources}
\begin{itemize}
\item[S1] ULAM.AI and the UnsolvedMath Contributors, supplied problems.json, record 2305052 (AMR-022-5052), AMR Open Problem Lists. Catalogue identity and source transcription; CC BY 4.0.
\url{https://huggingface.co/datasets/ulamai/UnsolvedMath}.
\item[S2] Hayman - Research Problems in Function Theory (2018), Problem 5.52. Original source indexed by AMR-022-5052.
\url{https://arxiv.org/abs/1809.07200}.
\end{itemize}
\end{document}
